\chapter{Sprint 2: Backend Services and API Integration}
\chaptertoc

\clearpage

\section{Introduction}
Sprint 2 expanded the project beyond a single local pipeline to support multiple camera sources. The backend was restructured to manage camera onboarding, feed control, and live delivery, providing a structured API that the frontend could depend on.

\section{Sprint Objectives and Backlog}
Sprint 2 transformed the backend from a single-pipeline script into a structured service layer capable of managing multiple camera sources, providing a reliable API for feed lifecycle operations, and delivering live streams with fallback awareness. The sprint objectives are captured in the user stories shown in Table~\ref{tab:sprint2_user_stories}, which outline manager needs for secure camera connectivity, automatic discovery, centralized feed management, and reliable live delivery.

\begin{table}[H]
    \centering
    \begingroup\setlength{\tabcolsep}{4pt}\renewcommand{\arraystretch}{0.95}\footnotesize
    \begin{tabularx}{\textwidth}{|l|X|l|X|}
        \hline
        User Story ID & User Story & Task ID & Task \\ \hline
        US2.1 & As a manager, I want IP camera streams connected securely so that live sources can be reused reliably. & T2.1 & Add authentication, transport control, reconnection, and snapshots for RTSP feeds. \\ \hline
        US2.2 & As a manager, I want cameras discovered automatically so that stream addresses do not need to be hardcoded. & T2.2 & Detect IP cameras and resolve their stream addresses. \\ \hline
        US2.3 & As a manager, I want feed management exposed through the API so that camera lifecycle actions are centralized. & T2.3 & Expose feed control, metadata handling, and diagnostic operations. \\ \hline
        US2.4 & As a manager, I want live delivery to stay available so that the dashboard can keep showing the active feed. & T2.4 & Support live delivery through real-time events and a layered video transport strategy. \\ \hline
    \end{tabularx}
    \caption{Sprint 2 user stories}
    \label{tab:sprint2_user_stories}
    \endgroup
\end{table}

\section{Design}

\subsection{Backend Domain Model and Class Structure}
The backend architecture organizes data entities hierarchically to support multi-manager and multi-feed scenarios while maintaining consistent state. The model structures manager accounts, feed configurations, video sessions, queue metrics, and alerts as a unified domain, where establishments parent both managers and feed configurations, and managers own feed definitions. This hierarchy ensures proper access scoping and state recovery across system restarts.

The core entities and their relationships are illustrated in Figure~\ref{fig:sprint2_class_diagram}.

\begin{figure}[!htbp]
    \centering
    \includegraphics[width=0.95\textwidth]{\detokenize{optional_class_diagram.pdf}}
    \caption{Sprint 2 backend class diagram showing core entities for feed orchestration and persistence.}
    \label{fig:sprint2_class_diagram}
\end{figure}

\section{Implementation}

\subsection{Source Integration and Camera Discovery}
The source integration layer supports multiple input types and network scenarios through a unified abstraction. IP cameras are connected via RTSP with secure credential handling, configurable TCP/UDP transport selection, connection validation, and automatic reconnection on stream interruption. Local files and webcams are routed through separate pathways to reduce complexity in the core source layer.

Camera discovery uses the ONVIF protocol to automatically identify cameras and retrieve stream information, eliminating hardcoded addresses. The implementation sends WS-Discovery probes over multicast UDP, processes ONVIF device responses, resolves media services, and requests RTSP URIs through SOAP calls. Streams are validated with OpenCV before registration to ensure connectivity. The workflow accommodates varying ONVIF implementations through a three-phase protocol: multicast UDP discovery, direct media service access for cameras that expose it, and device service indirection for cameras requiring capability queries first. This approach ensures robust discovery across heterogeneous camera deployments. The workflow was validated against a Digital Watchdog IP camera on a shared network.

The discovery process is exposed through backend endpoints for probe-based detection, stream resolution, camera testing, and snapshot capture. Credentials are handled securely without log exposure, and transport configuration adapts to stream requirements based on network conditions and source type.

The logical workflow of this discovery process is illustrated in Figure~\ref{fig:sprint2_camera_discovery_activity}.

\begin{figure}[!htbp]
    \centering
    \includegraphics[width=0.78\textwidth]{\detokenize{chapter4_camera_discovery_activity.pdf}}
    \caption{Camera discovery activity diagram: ONVIF discovery, RTSP resolution, and stream validation.}
    \label{fig:sprint2_camera_discovery_activity}
\end{figure}

The protocol-level details are captured in Figure~\ref{fig:chapter4_onvif_protocol_sequence}, which illustrates the three-phase ONVIF sequence: multicast UDP discovery, direct media service access, and device service indirection.

\begin{figure}[!htbp]
    \centering
    \includegraphics[width=0.95\textwidth]{\detokenize{chapter4_onvif_protocol_sequence.pdf}}
    \caption{ONVIF protocol sequence: WS-Discovery broadcast, direct media service resolution, and device service indirection.}
    \label{fig:chapter4_onvif_protocol_sequence}
\end{figure}


\subsection{API and Live Transport}
The backend exposes a structured REST and WebSocket API that coordinates feed management, source onboarding, metadata operations, system integration, and live event delivery. Access control is enforced at the session level: managers interact only with feeds in their account, while administrators maintain broader account governance scope. Access rules apply consistently across all operations including feed listing, creation, lifecycle control, snapshots, and real-time subscriptions.

Transport selection adapts to source type to minimize latency while maintaining fallback capability. IP camera feeds prioritize WebRTC for low-latency delivery, while local MP4 files use MJPEG endpoints. The backend communicates feed readiness, source type, and preferred playback paths to the frontend. For WebRTC-capable IP feeds, the backend facilitates browser offer/response negotiation through the media server. When the preferred transport is unavailable (feed down, service unreachable), the backend reports the condition explicitly, allowing the frontend to switch to fallback views rather than timing out silently. Transport configuration is externalized through environment variables for media URL, control URL, timeout, and feature flags.

Runtime control coordinates feed state, persistence synchronization, worker lifecycle management, and event dispatch through the live channel, which broadcasts snapshots, feed status changes, metric updates, alerts, and warnings. For local MP4 videos, a shared MJPEG stream ensures viewing continuity during runtime mode transitions.

\subsection{Persistence, State Recovery, and System Resilience}
The persistence layer uses SQLite to maintain establishments, managers, feed configurations, video sessions, alert history, and system state. The schema organizes entities hierarchically—establishments parent both managers and feed configurations; managers own feed definitions; sessions link back to establishments—ensuring proper access scoping and state isolation across restarts.

On startup, the persistence layer reloads saved feed definitions and resets any open sessions to a consistent state. This recovery mechanism prevents orphaned sessions and ensures feeds restore to their last known configuration without state inconsistency. The key architectural insight is that reliable connectivity extends beyond reaching a camera to maintaining application state across stream failures, transport mode switches, and system restarts. By explicitly managing recovery state and fallback conditions, the sprint ensured the backend could tolerate transient failures and resume operation reliably, a requirement critical for operational deployments.

\section{Conclusion}
Sprint 2 transformed the backend into a structured service layer capable of discovering and managing multiple camera sources, maintaining consistent feed state, and delivering live events reliably to the frontend. The sprint addressed three critical challenges: integrating heterogeneous camera types through ONVIF discovery and RTSP connectivity, maintaining access control and state consistency across multi-manager deployments, and implementing transport fallback strategies that ensure live delivery resilience.

The architectural achievement was recognizing that backend reliability requires more than connectivity—it demands explicit state recovery, transport flexibility, and graceful degradation when preferred transport becomes unavailable. By embedding these capabilities into the core backend design, Sprint 2 established a foundation that enabled straightforward frontend integration and positioned the system for operational deployment.


